\section{Motivation}
\label{sec:motivation}

\subsection{Limitations of Triple-Based Representations}
\label{sec:motivation:triples}

Standard knowledge graphs decompose all facts into subject-predicate-object
triples. While this uniform representation simplifies storage and querying
for binary relations, it introduces significant overhead for $n$-ary
relations such as ``Alice sold a car to Bob for \$20,000 on 2024-01-15.''
This fact has arity $k = 4$ (buyer, seller, object, price, date) and
requires reification into $k + 2$ triples~\cite{manola2004rdf}.

The consequences include:
\begin{itemize}
    \item \textbf{Join explosion:} Reconstructing the original fact requires
          $k$ joins, each with cost proportional to the index size.
    \item \textbf{Storage amplification:} Each $n$-ary fact generates
          $n+3$ triples instead of a single record.
    \item \textbf{Loss of context:} The context of a fact (its scope of
          validity) is not a first-class concept in the triple model.
\end{itemize}

\subsection{The Case for Sheaf Theory}
\label{sec:motivation:sheaf}

Sheaf theory addresses these limitations by providing:
\begin{itemize}
    \item \textbf{First-class $n$-ary relations:} A fact with arity $k$
          is stored as a single section, not decomposed.
    \item \textbf{Native context support:} The poset of contexts models
          scoping directly.
    \item \textbf{Locality:} Queries within a single context examine only
          that context's sections.
    \item \textbf{Gluing:} Global facts can be reconstructed from compatible
          local sections.
\end{itemize}

\subsection{Research Gap}
\label{sec:motivation:gap}

While sheaf theory has been applied to distributed computing, data fusion,
and sensor networks~\cite{robinson2020sheaf}, no prior work has built a
complete database system around sheaf-theoretic principles and evaluated it
against a traditional KG baseline on realistic workloads.
